% chapter4_future_work.tex
\chapter{Research Gaps \& Future Work}
\label{ch:future-work}

While this study provides a comprehensive empirical foundation for understanding LLM citation behavior, several avenues remain open for future investigation. The following subsections outline the most pressing research gaps identified during this work.

\subsection{Deeper Crawling}
\label{sec:fw-deeper-crawling}

We stopped at Rank 200; going to 300+ might find more matches. Expanding the crawl depth would help determine whether additional ``invisible'' citations can be recovered at greater SERP depths, or whether a hard ceiling exists beyond which search-index overlap ceases to grow.

\subsection{Longitudinal Consistency (Expanded Runs)}
\label{sec:fw-longitudinal-consistency}

While this study used 3--4 runs per prompt, future work should expand this to 10+ runs to achieve statistical significance in ``stochastic retrieval'' patterns and to better map the ``long tail'' of citations that appear only in rare instances.

\subsection{Temporal Analysis}
\label{sec:fw-temporal-analysis}

How do results change over time? Running the same queries in 3 months would reveal whether citation patterns are stable or subject to temporal drift as model weights, retrieval indices, and web content evolve.

\subsection{Query Category Segmentation}
\label{sec:fw-query-category}

Do certain product categories have better or worse overlap? Segmenting the dataset by category (e.g., electronics vs.\ fashion vs.\ software) could reveal category-specific citation dynamics that are masked in aggregate analysis.

\subsection{Multi-Model Comparison}
\label{sec:fw-multi-model}

Compare ChatGPT vs.\ Gemini vs.\ Claude on the same queries. While this study includes both ChatGPT and Gemini, extending the comparison to additional frontier models would strengthen the generalisability of the findings.

\subsection{User Study}
\label{sec:fw-user-study}

Do humans prefer ChatGPT's recommendations or Bing's Top~10? A controlled user study measuring preference, trust, and purchase intent would complement the purely structural analysis presented here.

\subsection{Exhaustive SERP Depth Analysis}
\label{sec:fw-serp-depth}

Expanding Bing/Google search depth from top 200 to top 300+ results would help determine if ``invisible'' citations are simply lower-ranked search results or truly independent LLM retrievals.

\subsection{Residual Source Isolation}
\label{sec:fw-residual-source}

By programmatically ``subtracting'' all possible SERP matches (even at extreme depths), researchers can isolate the true ``LLM-native'' grounding set---sources ChatGPT/Gemini access via internal knowledge bases, direct partnerships, or non-public indices.

\subsection{Decay Rate of Attribution}
\label{sec:fw-decay-rate}

Analyzing if the probability of an LLM citing a source correlates with its SERP rank even beyond the first few pages, or if the LLM's ``internal'' prioritization overrides search engine ranking at depth.
